\documentclass[10pt,a4paper]{amsart}
\usepackage[margin=19mm]{geometry}
\usepackage{amsmath,amssymb,mathtools}
\usepackage[scr=rsfs,cal=euler]{mathalfa}
\usepackage{xfrac}
\usepackage[unicode,hidelinks,bookmarks=false]{hyperref}
\newcommand{\ie}{i.e.\ }
\newcommand{\cf}{cf.\ }
\newcommand{\resp}{respectively\ }
\newcommand{\Subsection}{\subsection}
\providecommand{\nobreakdash}{\nobreak\hbox{-}}
\setlength{\emergencystretch}{2em}
\multlinegap0pt
\usepackage{amssymb}
\newtheorem{lemma}{Lemma}
\title{Convergence to the equilibrium for the kinetic transport equation in the two-dimensional periodic Lorentz Gas}
\author{Francesca Pieroni}
\date{Source: arXiv:2504.06839v3 (CC BY 4.0)}
\begin{document}
\maketitle

\noindent\emph{Source and licence.} Convergence to the equilibrium for the kinetic transport equation in the two-dimensional periodic Lorentz Gas. Version arXiv:2504.06839v3 (CC BY 4.0), distributed by arXiv under the Creative Commons Attribution 4.0 International licence, http://creativecommons.org/licenses/by/4.0/, as recorded in the arXiv licence metadata for this exact version and shown on the abstract page https://arxiv.org/abs/2504.06839v3. Full licence notice: \texttt{licenses/arxiv-2504.06839-CC-BY-4.0.txt}.\par\smallskip

\noindent\textit{Editorial scope.} Counted unit: the article's lemma lemma:rho (source0.tex lines 1192--1224). The article's complete original proof of it is delivered with it (source0.tex lines 1579--1754, one \texttt{proof} environment of 176 lines, which the article places away from the statement). No mathematics is written, repaired or re-derived: the statement and the proof are the article's own lines, in the article's own notation, and no retained line is substituted or re-flowed.  Retained uncounted as the source-local context the proof needs: lines 199--201; lines 205--209; lines 1134--1136; lines 1140--1144; lines 1148--1152; lines 1156--1160; lines 2482--2488.  Nothing was omitted from any retained span, so the package carries no omission record.  No retained line contains a citation, so the wrapper carries no bibliography and no bibliography entry.  No retained line contains a figure environment, an image-inclusion command or drawing code, so no image is delivered and the package declares no asset dependency.  Editorial formatting only, in addition: an AMS \texttt{amsart} preamble that restates the article's own declarations and macros used by the retained lines, namely the environments \texttt{lemma} on separate counters, together with the standard packages those lines need.  The retained environments print in this document as Lemma 1 in order of appearance, whereas the article prints the same items with the numbering of its own full document.  The complete native source of the article as distributed in the arXiv e-print for this exact version is bundled unchanged as \texttt{sources/176243/source0.tex}.
{\tiny\begin{align}\label{eq:ev}
\left\{\begin{array}{lcr}\frac{\partial\mu_t}{\partial t}(x,\theta,s,h)+v(\theta)\cdot\nabla_x\mu_t(x,\theta,s,h)-\frac{\partial\mu_t}{\partial s}(x,\theta,s,h)=\int_{-1}^1dh' Q(s,h|h')\mu_t(x,\theta+\pi-2\arcsin(h'),0,h'),\\\mu_t(x,\theta,s,h)|_{t=0}=\mu_0(x,\theta,s,h).\end{array}\right.
\end{align}}

{\footnotesize\begin{align}
\nonumber\mu_t(x,\theta,s,h)&=\mu_0(x-tv(\theta),\theta,s+t,h)
\\
\label{rapp:ev}&+\int_0^tdt'\int_{-1}^1dh'Q(s+t-t',h|h')\mu_{t'}(x-(t-t')v(\theta),\theta+\pi-2\arcsin(h'),0,h'),
\end{align}}

{\small\begin{align}\label{mutL1cpt}
\int_{\mathbb{X}}dx\int_{\mathbb{T}^1_{2\pi}}d\theta\int_0^Tds\int_{-1}^1dh|\mu_t(x,\theta,s,h)|\leq\int_{\mathbb{X}}dx\int_{\mathbb{T}^1_{2\pi}}d\theta\int_0^{T+t}ds\int_{-1}^1dh|\mu_0(x,\theta,s,h)|,
\end{align}}

{\small\begin{align}
\label{distanzadecr}\int_{\mathbb{X}}dx\int_{\mathbb{T}^1_{2\pi}}d\theta\int_0^{\infty}ds\int_{-1}^1dh|\mu_t(x,\theta,s,h)|&\leq\int_{\mathbb{X}}dx\int_{\mathbb{T}^1_{2\pi}}d\theta\int_0^{\infty}ds\int_{-1}^1dh|\mu_0(x,\theta,s,h)|,
\\
\label{consmassa}\int_{\mathbb{X}}dx\int_{\mathbb{T}^1_{2\pi}}d\theta\int_0^{\infty}ds\int_{-1}^1dh\mu_t(x,\theta,s,h)&=\int_{\mathbb{X}}dx\int_{\mathbb{T}^1_{2\pi}}d\theta\int_0^{\infty}ds\int_{-1}^1dh\mu_0(x,\theta,s,h).
\end{align}}

\begin{align}
\nonumber\|\mu_t\|_{L^p(\mathbb{X}\times\mathbb{T}^1_{2\pi}\times[0,T]\times[-1,1])}&\leq C\|\mu_0\|_{L^p(\mathbb{R}^2\times\mathbb{T}^1_{2\pi}\times[0,t+T]\times[-1,1])}
\\
\label{normaLp}&+C\int_{\mathbb{T}^1_{2\pi}}d\theta'\int_0^tdt'\int_{-1}^1dh'\|\mu_0(\cdot,\theta',t',h')\|_{L^p(\mathbb{X})},
\end{align}

\begin{align}
\nonumber\|\mu_t\|_{L^p(\mathbb{X}\times\mathbb{T}^1_{2\pi}\times[0,+\infty)\times[-1,1])}&\leq C\|\mu_0\|_{L^p(\mathbb{R}^2\times\mathbb{T}^1_{2\pi}\times[0,+\infty)\times[-1,1])}
\\
\label{normaLpfinita}&+C\int_{\mathbb{T}^1_{2\pi}}d\theta'\int_0^tdt'\int_{-1}^1dh'\|\mu_0(\cdot,\theta',t',h')\|_{L^p(\mathbb{X})}.
\end{align}

{\small\begin{align}
\nonumber&\left|\int_{t_1}^{t_2}dt'\int_{-1}^1dh'Q(t-t',h|h')\mu_0(\theta,t',h')\right|
\\
\nonumber&\leq\int_{t_1}^{t_2}dt'\int_{-1}^1dh'Q(t-t',h|h')|\mu_0(\theta,t',h')|^p\left(\int_{t_1}^{t_2}dt'\int_{-1}^1dh'Q(t-t',h|h')\right)^{p-1}
\\
\label{Jensencons}&=\int_{t_1}^{t_2}dt'\int_{-1}^1dh'Q(t-t',h|h')|\mu_0(\theta,t',h')|^p\left(E(t-t_2,h)-E(t-t_1,h)\right)^{p-1}.
\end{align}}

\begin{lemma}\label{lemma:rho}
Let $\mathbb{X}=\mathbb{R}^2$ or $\mathbb{X}=\mathbb{T}^2=\mathbb{R}^2/\mathbb{Z}^2$, $T>0$, $p\in[1,+\infty]$ and $\mu\in L^p(\mathbb{R}^2\times\mathbb{T}^1_{2\pi}\times[0,T]\times[-1,1])$. Then there exists a unique $\rho\in L^p(\mathbb{R}^2\times\mathbb{T}^1_{2\pi}\times[0,T]\times[-1,1])$ such that

{\small\begin{align}
\nonumber\rho(x,\theta,t,h)&=\mu(x-tv(\theta),\theta,t,h)
\\
\label{proprho}&+\int_0^tdt'\int_{-1}^1dh'Q(t-t',h|h')\rho(x-(t-t')v(\theta),\theta+\pi-2\arcsin(h'),t',h').
\end{align}}

Moreover for any $t\leq T$ we have

{\footnotesize\begin{align}
\label{propr1}&\int_{\mathbb{X}}dx\int_{\mathbb{T}^1_{2\pi}}d\theta \int_0^tdt'\int_{-1}^1dh|\rho(x,\theta,t',h)|E(t-t',h)\leq \int_{\mathbb{X}}dx\int_{\mathbb{T}^1_{2\pi}}d\theta \int_0^tdt'\int_{-1}^1dh|\mu(x,\theta,t',h)|,
\\
\label{propr2}&\int_{\mathbb{X}}dx\int_{\mathbb{T}^1_{2\pi}}d\theta \int_0^tdt'\int_{-1}^1dh\rho(x,\theta,t',h)E(t-t',h)=\int_{\mathbb{X}}dx\int_{\mathbb{T}^1_{2\pi}}d\theta \int_0^tdt'\int_{-1}^1dh\mu(x,\theta,t',h).
\end{align}}

Moreover, there exists a constant $C>0$ such that

\begin{align}
\nonumber\|\rho(\cdot,\theta,t,h)\|_{L^p(\mathbb{X})}&\leq\|\mu(\cdot,\theta,t,h)\|_{L^p(\mathbb{X})}
\\
\nonumber&+\int_0^tdt'\int_{-1}^1dh'Q(t-t',h|h')\|\mu(\cdot,\theta+\pi-2\arcsin(h'),t',h')\|_{L^p(\mathbb{X})}
\\
\label{normaLprho}&+C\int_{\mathbb{T}^1_{2\pi}}d\theta'\int_0^tdt'\int_{-1}^1dh'\|\mu(\cdot,\theta',t',h')\|_{L^p(\mathbb{X})}.
\end{align}

Furthermore, the following properties hold:
\begin{itemize}
\item if $\mu\geq0$ then $\rho\geq0$,
\item if $\mu$ does not depend on $x$ or $(x,\theta)$ neither does $\rho$. In such case, \eqref{propr1} and \eqref{propr2} hold also without integrating with respect to $x$ or $(x,\theta)$.
\end{itemize}
\end{lemma}

\begin{proof}
We start by proving the existence and the uniqueness: Lemma \ref{lemma:rho} with $\mathbb{X}=\mathbb{R}^2$ or $\mathbb{X}=\mathbb{T}^2$, $p\in[1,+\infty]$ and $\mu=\mu_0$ provides the existence of a function $\rho$ satisfying \eqref{proprho}. Therefore, if we define $\mu_t(x,\theta,0,h)=\rho(x,\theta,t,h)$, and we then use equation \eqref{rapp:ev} to obtain $\mu_t(x,\theta,s,h)$, we get a mild solution of the equation \eqref{eq:ev}. The uniqueness is a consequence of the fact that the relation \eqref{rapp:ev} forces $\mu_t(x,\theta,0,h)$ to be a solution of \eqref{proprho}, which is unique thanks to Lemma \ref{lemma:rho}, and therefore also $\mu_t(x,\theta,s,h)$ is uniquely defined for any $s\geq0$.

The next step is to prove property \eqref{mutL1cpt}.

By representation formula \eqref{rapp:ev} we get

{\tiny\begin{align*}
&\int_{\mathbb{X}}dx\int_{\mathbb{T}^1_{2\pi}}d\theta\int_0^Tds\int_{-1}^1dh|\mu_t(x,\theta,s,h)|
\\
&\leq\int_{\mathbb{X}}dx\int_{\mathbb{T}^1_{2\pi}}d\theta\int_0^Tds\int_{-1}^1dh|\mu_0(x-tv(\theta),\theta,s+t,h)|
\\
&+\int_{\mathbb{X}}dx\int_{\mathbb{T}^1_{2\pi}}d\theta\int_0^Tds\int_{-1}^1dh\int_0^tdt'\int_{-1}^1dh' Q(s+t-t',h|h')|\mu_{t'}(x-(t-t')v(\theta),\theta+\pi-2\arcsin(h'),0,h')|
\\
&=\int_{\mathbb{X}}dx\int_{\mathbb{T}^1_{2\pi}}d\theta\int_t^{t+T}ds\int_{-1}^1dh|\mu_0(x,\theta,s,h)|
\\
&+\int_{\mathbb{X}}dx\int_{\mathbb{T}^1_{2\pi}}d\theta\int_0^Tds\int_{-1}^1dh\int_0^tdt'\int_{-1}^1dh' Q(s+t-t',h|h')|\mu_{t'}(x,\theta,0,h')|
\\
&=\int_{\mathbb{X}}dx\int_{\mathbb{T}^1_{2\pi}}d\theta\int_t^{t+T}ds\int_{-1}^1dh|\mu_0(x,\theta,s,h)|
\\
&+\int_{\mathbb{X}}dx\int_{\mathbb{T}^1_{2\pi}}d\theta\int_0^tdt'\int_{-1}^1dh' |\mu_{t'}(x,\theta,0,h')|\int_0^Tds\int_{-1}^1dhQ(s+t-t',h|h').
\end{align*}}

Now we look at the last term, and since
\[
\int_0^Tds\int_{-1}^1dhQ(s+t-t',h|h')=E(t-t',h')-E(t+T-t',h')\leq E(t-t',h'),
\]
we can write

\begin{align*}
&\int_{\mathbb{X}}dx\int_{\mathbb{T}^1_{2\pi}}d\theta\int_0^Tds\int_{-1}^1dh|\mu_t(x,\theta,s,h)|
\\
&\leq\int_{\mathbb{X}}dx\int_{\mathbb{T}^1_{2\pi}}d\theta\int_t^{t+T}ds\int_{-1}^1dh|\mu_0(x,\theta,s,h)|
\\
&+\underbrace{\int_{\mathbb{X}}dx\int_{\mathbb{T}^1_{2\pi}}d\theta\int_0^tdt'\int_{-1}^1dh' |\mu_{t'}(x,\theta,0,h')|E(t-t',h')}_{\leq\int_{\mathbb{X}}dx\int_{\mathbb{T}^1_{2\pi}}d\theta\int_0^tdt'\int_{-1}^1dh' |\mu_0(x,\theta,t',h')|\textrm{ by }\eqref{propr1}\textrm{ of Lemma \ref{lemma:rho}}}
\\
&\leq\int_{\mathbb{X}}dx\int_{\mathbb{T}^1_{2\pi}}d\theta\int_t^{t+T}ds\int_{-1}^1dh|\mu_0(x,\theta,s,h)|
\\
&+\int_{\mathbb{X}}dx\int_{\mathbb{T}^1_{2\pi}}d\theta\int_0^tdt'\int_{-1}^1dh' |\mu_0(x,\theta,t',h')|
\\
&=\int_{\mathbb{X}}dx\int_{\mathbb{T}^1_{2\pi}}d\theta\int_0^{t+T}ds\int_{-1}^1dh|\mu_0(x,\theta,s,h)|,
\end{align*}

and this proves property \eqref{mutL1cpt}.

As for \eqref{distanzadecr}, it is sufficient to take the limit $T\to+\infty$ in the equation \eqref{mutL1cpt}.

Then, to prove property \eqref{consmassa}, the same argument as the one used to prove property \eqref{mutL1cpt} works: the only differences are that \eqref{propr2} should be used instead of  \eqref{propr1} and that all the computations have to be meant in the limit $T\to+\infty$.

The positivity of $\mu_t$ is preserved in time since Lemma \ref{lemma:rho} ensures that $\mu_0\geq0$ a.e. implies $\rho(x,\theta,t,h)=\mu_t(x,\theta,0,h)\geq0$ a.e. and therefore thanks to \eqref{rapp:ev} we have also $\mu_t\geq0$ a.e..

Moreover, if $\mu_0$ does not depend on $x$ or on $(x,\theta)$ neither $\rho(x,\theta,t,h)=\mu_t(x,\theta,0,h)$ does thanks to Lemma \ref{lemma:rho}. Therefore, thanks to the representation formula \eqref{rapp:ev}, the same holds for $\mu_t$.

Now we only have to prove \eqref{normaLp}. To this purpose, if we use the representation \eqref{rapp:ev}, with triangle inequality and by extending the $L^p$ norm of the second summand to all $s\in[0,+\infty)$, we get

{\scriptsize\begin{align}
\label{normaLpstep1}\|\mu_t\|_{L^p(\mathbb{X}\times\mathbb{T}^1_{2\pi}\times[0,T]\times[-1,1])}&\leq\|\mu_0(\cdot,\cdot,\cdot+t,\cdot)\|_{L^p(\mathbb{X}\times\mathbb{T}^1_{2\pi}\times[0,T]\times[-1,1])}
\\
\label{normaLpstep2}&+\left\|\int_0^tdt'\int_{-1}^1dh'Q(\cdot+t-t',\cdot|h')\mu_{t'}(\cdot,\cdot,0,h')\right\|_{L^p(\mathbb{X}\times\mathbb{T}^1_{2\pi}\times[0,+\infty)\times[-1,1])}.
\end{align}}

Now since the first term in \eqref{normaLpstep1} can be expressed as
\[
\eqref{normaLpstep1}=\|\mu_0(\cdot,\cdot,\cdot+t,\cdot)\|_{L^p(\mathbb{X}\times\mathbb{T}^1_{2\pi}\times[0,T]\times[-1,1])}=\|\mu_0\|_{L^p(\mathbb{X}\times\mathbb{T}^1_{2\pi}\times[t,T+t]\times[-1,1])},
\]
we focus on the second summand. Since
\[
\|\cdot\|_{L^p(\mathbb{X}\times\mathbb{T}^1_{2\pi}\times[0,+\infty)\times[-1,1])}=\left\|\|\cdot\|_{L^p(\mathbb{X})}\right\|_{L^p(\mathbb{T}^1_{2\pi}\times[0,+\infty)\times[-1,1])},
\]
using again Minkowski's inequality, we can write

\begin{align*}
\eqref{normaLpstep2}&=\left\|\int_0^tdt'\int_{-1}^1dh'Q(\cdot+t-t',\cdot|h')\mu_{t'}(\cdot,\cdot,0,h')\right\|_{L^p(\mathbb{X}\times\mathbb{T}^1_{2\pi}\times[0,+\infty)\times[-1,1])}
\\
&\leq\left\|\int_0^tdt'\int_{-1}^1dh'Q(\cdot+t-t',\cdot|h')\|\mu_{t'}(\cdot,\cdot,0,h')\|_{L^p(\mathbb{X})}\right\|_{L^p(\mathbb{T}^1_{2\pi}\times[0,+\infty)\times[-1,1])},
\end{align*}

and using \eqref{normaLprho} of Lemma \ref{lemma:rho} and applying triangle inequality to the external norm $\|\|_{L^p(\mathbb{T}^1_{2\pi}\times[0,+\infty)\times[-1,1])}$, we have

\begin{align*}
\eqref{normaLpstep2}&\leq\left\|\int_0^tdt'\int_{-1}^1dh'Q(\cdot+t-t',\cdot|h')\left[\|\mu_0(\cdot,\cdot,t',h')\|_{L^p(\mathbb{X})}\right.\right.
\\
&+\int_0^{t'}dt''\int_{-1}^1dh''Q(t'-t'',h'|h'')\|\mu_0(\cdot,\cdot+\pi-2\arcsin(h''),t'',h'')\|_{L^p(\mathbb{X})}
\\
&+C\left.\left.\int_{\mathbb{T}^1_{2\pi}}d\theta'\int_0^{t'}dt''\int_{-1}^1dh''\|\mu_0(\cdot,\theta',t'',h'')\|_{L^p(\mathbb{X})}\right]\right\|_{L^p(\mathbb{T}^1_{2\pi}\times[0,+\infty)\times[-1,1])}.
\end{align*}

By triangle inequality we get

{\footnotesize\begin{align}
\nonumber&\eqref{normaLpstep2}
\\
\label{normaLpstep3}&\leq\left\|\int_0^tdt'\int_{-1}^1dh'Q(\cdot+t-t',\cdot|h')\|\mu_0(\cdot,\cdot,t',h')\|_{L^p(\mathbb{X})}\right\|_{L^p}
\\
\label{normaLpstep5}&+\left\|\int_0^tdt'\int_{-1}^1dh'Q(\cdot+t-t',\cdot|h')\int_0^{t'}dt''\int_{-1}^1dh'' Q(t'-t'',h'|h'')\|\mu_0(\cdot,\theta',t'',h'')\|_{L^p(\mathbb{X})}\right\|_{L^p}
\\
\label{normaLpstep4}&+C\left\|\int_0^tdt'\int_{-1}^1dh'Q(\cdot+t-t',\cdot|h')\int_0^{t'}dt''\int_{-1}^1dh''\|\mu_0(\cdot,\cdot,t'',h'')\|_{L^p(\mathbb{X})}\right\|_{L^p},
\end{align}}

where in the previous inequality we shortened
\[
L^p:=L^p(\mathbb{T}^1_{2\pi}\times[0,+\infty)\times[-1,1]),
\]
and
\[
\theta':=\theta+\pi-2\arcsin(h'').
\]
Hereafter, we will assume that $p$ is finite, but the case $p=\infty$ can be studied in the same way.

By using Jensen's inequality in the inner integral as in \eqref{Jensencons} and changing the integration order, one has

{\tiny\begin{align*}
&\eqref{normaLpstep3}
\\
&\leq\left(\int_{\mathbb{T}^1_{2\pi}}d\theta\int_0^{\infty}ds\int_{-1}^1dh\int_0^tdt'\int_{-1}^1dh'Q(s+t-t',h|h')\|\mu_0(\cdot,\theta,t',h')\|_{L^p(\mathbb{X})}^p\underbrace{(E(s,h)-E(s+t,h))^{p-1}}_{\leq1}\right)^{\frac{1}{p}}
\\
&=\left(\int_{\mathbb{T}^1_{2\pi}}d\theta\int_0^tdt'\int_{-1}^1dh'\|\mu_0(\cdot,\theta,t',h')\|_{L^p(\mathbb{X})}^p\underbrace{\int_0^{\infty}ds\int_{-1}^1dhQ(s+t-t',h|h')}_{=E(t-t',h')\leq1}\right)^{\frac{1}{p}}
\\
&=\|\mu_0\|_{L^p(\mathbb{X}\times\mathbb{T}^1_{2\pi}\times[0,t]\times[-1,1])},
\end{align*}}

and, first applying again twice Jensen's inequality as in \eqref{Jensencons}, and then changing the integration order, one has also

{\tiny\begin{align*}
\eqref{normaLpstep5}^p\leq\int d\theta dsdh\int_0^tdt'\int_{-1}^1dh'Q(s+t-t',h|h')\int_0^{t'}dt''\int_{-1}^1dh''Q(t'-t'',h'|h'')\|\mu_0(\cdot,\theta,t'',h'')\|_{L^p(\mathbb{X})}^p\xi,
%&=\left\|\int_0^tdt'\int_{-1}^1dh'Q(\cdot+t-t',\cdot|h') \int_0^{t'}dt''\int_{-1}^1dh''Q(t'-t'',h'|h'')\|\mu(\cdot,\cdot,t'',h'')\|_{L^p(\mathbb{X})}\right\|_{L^p(\mathbb{T}^1_{2\pi}\times[0,+\infty)\times[-1,1])}
%\\
\end{align*}}

with
\[
\xi=\xi(t,s,h,t',h'):=\underbrace{(1-E(t',h'))^{p-1}}_{\leq1}\underbrace{(E(s,h)-E(s+t,h))^{p-1}}_{\leq 1})^{\frac{1}{p}}\leq1,
\]
and
\[
\int d\theta dsdh:=\int_{\mathbb{T}^1_{2\pi}}d\theta\int_0^{\infty}ds\int_{-1}^1dh.
\]
Therefore

{\small\begin{align*}
\eqref{normaLpstep5}\leq\left(\int_{\mathbb{T}^1_{2\pi}}d\theta\int_0^tdt''\int_{-1}^1dh''\|\mu_0(\cdot,\theta,t'',h'')\|_{L^p(\mathbb{X})}^p\int_{t''}^tdt'\int_{-1}^1dh'Q(t'-t'',h'|h'')\xi_1\right)^{\frac{1}{p}}
\end{align*}}

with
\[
\xi_1=\xi_1(t,t',h'):=\int_0^{\infty}ds\int_{-1}^1dhQ(s+t-t',h|h')=E(t-t',h')\leq1.
\]
Hence

{\footnotesize\begin{align*}
\eqref{normaLpstep5}&\leq\left(\int_{\mathbb{T}^1_{2\pi}}d\theta\int_0^tdt''\int_{-1}^1dh''\|\mu_0(\cdot,\theta,t'',h'')\|_{L^p(\mathbb{X})}^p\underbrace{\int_{t''}^tdt'\int_{-1}^1dh'Q(t'-t'',h'|h'')}_{=E(t-t'',h'')\leq1}\right)^{\frac{1}{p}}
\\
&\leq\left(\int_{\mathbb{T}^1_{2\pi}}d\theta\int_0^tdt''\int_{-1}^1dh''\|\mu_0(\cdot,\theta,t'',h'')\|_{L^p(\mathbb{X})}^p\right)^{\frac{1}{p}}
\\
&=\|\mu_0\|_{L^p(\mathbb{X}\times\mathbb{T}^1_{2\pi}\times[0,t]\times\mathbb[-1,1])}.
\end{align*}}

Finally, one has

{\footnotesize\begin{align*}
\eqref{normaLpstep4}&=C\left\|\int_0^tdt'\int_{-1}^1dh'Q(\cdot+t-t',\cdot|h')\underbrace{\int_{\mathbb{T}^1_{2\pi}}d\theta'\int_0^{t'}dt''\int_{-1}^1dh''\|\mu_0(\cdot,\theta',t'',h'')\|_{L^p(\mathbb{X})}}_{\leq\int_{\mathbb{T}^1_{2\pi}}d\theta'\int_0^tdt''\int_{-1}^1dh''\|\mu_0(\cdot,\theta',t'',h'')\|_{L^p(\mathbb{X})}}\right\|_{L^p}
\\
&\leq C\int_{\mathbb{T}^1_{2\pi}}d\theta'\int_0^tdt''\int_{-1}^1dh''\|\mu_0(\cdot,\theta',t'',h'')\|_{L^p(\mathbb{X})}\left\|\underbrace{\int_0^tdt'\int_{-1}^1dh'Q(\cdot+t-t',\cdot|h')}_{=E-E(\cdot+t,\cdot)\leq E}\right\|_{L^p}
\\
&\leq(2\pi)^{\frac{1}{p}} C\int_{\mathbb{T}^1_{2\pi}}d\theta'\int_0^tdt''\int_{-1}^1dh''\|\mu_0(\cdot,\theta',t'',h'')\|_{L^p(\mathbb{X})},
\end{align*}}

where, in the first two lines of the previous inequality, we shortened
\[
L^p:=L^p(\mathbb{T}^1_{2\pi}\times[0,+\infty)\times[-1,1]).
\]

Summing \eqref{normaLpstep1} and \eqref{normaLpstep2}, which is smaller than \eqref{normaLpstep3}+\eqref{normaLpstep5}+\eqref{normaLpstep4}, we get \eqref{normaLp}.

Finally, \eqref{normaLpfinita} is obtained by sending $T\to+\infty$ in \eqref{normaLp}.
\end{proof}

\end{document}
